% Generated draft: requires mathematical and source review.
\documentclass[11pt]{article}
\usepackage[margin=25mm]{geometry}
\usepackage{fontspec}
\setmainfont{Latin Modern Roman}
\usepackage{amsmath,amssymb,mathtools,unicode-math}
\setmathfont{Latin Modern Math}
\usepackage{enumitem,needspace,xurl,hyperref}
\setlength{\parindent}{0pt}
\setlength{\parskip}{5pt}
\setlength{\emergencystretch}{3em}
\newcommand{\N}{\mathbb{N}}
\newcommand{\Z}{\mathbb{Z}}
\newcommand{\Q}{\mathbb{Q}}
\newcommand{\R}{\mathbb{R}}
\newcommand{\C}{\mathbb{C}}
\newcommand{\F}{\mathbb{F}}
\newcommand{\A}{\mathbb{A}}
\newcommand{\PP}{\mathbb{P}}
\newcommand{\E}{\mathbb{E}}
\newcommand{\eps}{\varepsilon}
\newcommand{\dd}{\,\mathrm{d}}
\begin{document}
\section*{UnsolvedMath ID: 800595}
\textbf{Source problem code:} catalogue entry
\subsection{Definitions and mathematical statement}
The supplied record does not specify the underlying class of spaces or define the cover classes. The following gives an explicit formulation under the usual countable Borel-cover conventions; these conventions and the intended domain require confirmation against the source.

Let \(X\) be a topological space. A countable Borel cover of \(X\) is an at most countable family \(\mathcal U\) of Borel subsets of \(X\) such that \(\bigcup\mathcal U=X\) and \(X\notin\mathcal U\). Write \(\mathrm B(X)\) for the class of these covers. Define the following subclasses:
\begin{itemize}
\item \(\mathrm B_\Omega(X)\) consists of those \(\mathcal U\in\mathrm B(X)\) such that every finite subset of \(X\) is contained in some member of \(\mathcal U\).
\item \(\mathrm B_\Gamma(X)\) consists of those infinite \(\mathcal U\in\mathrm B(X)\) such that, for every \(x\in X\), the family \(\{U\in\mathcal U:x\notin U\}\) is finite.
\item \(\mathrm B_{\mathrm T}(X)\) consists of those \(\mathcal U\in\mathrm B(X)\) such that every point of \(X\) belongs to infinitely many members of \(\mathcal U\), and, for every \(x,y\in X\), at least one of the families
\[
\{U\in\mathcal U:x\in U,\ y\notin U\},\qquad
\{U\in\mathcal U:y\in U,\ x\notin U\}
\]
is finite. These are the Borel \(\tau\)-covers.
\end{itemize}
For classes \(\mathcal A(X)\) and \(\mathcal D(X)\) of covers, the property \(\mathsf S_1(\mathcal A,\mathcal D)\) means that, for every sequence \((\mathcal U_n)_{n\geq1}\) from \(\mathcal A(X)\), there are choices \(U_n\in\mathcal U_n\) such that \(\{U_n:n\geq1\}\in\mathcal D(X)\). The property \(\mathsf S_{\mathrm{fin}}(\mathcal A,\mathcal D)\) means that, for every such sequence, there are finite subfamilies \(\mathcal F_n\subseteq\mathcal U_n\) such that \(\bigcup_{n\geq1}\mathcal F_n\in\mathcal D(X)\).

For an intended class \(\mathcal C\) of spaces, call a property \(P\) hereditary on \(\mathcal C\) if
\[
\forall X\in\mathcal C\ \forall Y\subseteq X,
\qquad P(X)\Longrightarrow P(Y),
\]
where \(Y\) has the subspace topology and its cover classes are formed using Borel subsets of \(Y\).

The question is which, if any, of the following six properties are hereditary:
\[
\begin{gathered}
\mathsf S_1(\mathrm B_{\mathrm T},\mathrm B_\Gamma),\qquad
\mathsf S_1(\mathrm B_{\mathrm T},\mathrm B_{\mathrm T}),\qquad
\mathsf S_1(\mathrm B_{\mathrm T},\mathrm B_\Omega),\qquad
\mathsf S_1(\mathrm B_{\mathrm T},\mathrm B),\\
\mathsf S_{\mathrm{fin}}(\mathrm B_{\mathrm T},\mathrm B_{\mathrm T}),\qquad
\mathsf S_{\mathrm{fin}}(\mathrm B_{\mathrm T},\mathrm B_\Omega).
\end{gathered}
\]
The class \(\mathcal C\) is unresolved in the supplied record. The cited title suggests a setting involving sets of reals, but the record does not establish the exact ambient-space hypotheses. Thus this is a conditional formulation pending confirmation of the domain and cover conventions.
\subsection{Short English statement}
For each of six selection principles involving countable Borel \(\tau\)-covers, ask whether every subspace of a space satisfying the principle also satisfies it. The supplied record leaves the intended class of spaces and the precise cover conventions unstated.
\subsection{Sources}
\begin{itemize}
\item Local UnsolvedMath registry, permanent ID \(800595\), source problem code \texttt{catalogue entry}; record title: ``Selection Principles in Mathematics --- Problem 23 ([ideals]).'' The supplied record attributes the problem to Tsaban, gives the proposed year as 2006, and labels its status ``open.'' This status is supplied registry metadata, not an independently verified assessment.
\item Supplied source list: Tsaban, \emph{Selection Principles and Special Sets of Reals: Open Problems} (2006), Problem 23. Supplied URL: \url{https://arxiv.org/abs/math/0609601}. The record reports access on 2026-10-16 and extraction by \texttt{arxiv\_latex}; these are supplied provenance details, not a claim of source access or verification for this draft.
\item The supplied attribution/context annotation is \texttt{[ideals]}. Its expansion and any associated reference are absent from the record. No additional literature references or dated literature assessment are supplied.
\end{itemize}
\end{document}
